\documentclass{book}

% Pacchetti di Codifica e Lingua
\usepackage[utf8]{inputenc}   % Codifica dei caratteri
\usepackage[T1]{fontenc}      % Codifica dei font
\usepackage[italian]{babel}   % Lingua italiana

% Pacchetti Grafici e Diagrammi
\usepackage{graphicx}         % Inserimento immagini
\usepackage{tikz}             % Rappresentazioni grafiche
\usetikzlibrary{angles,quotes,arrows.meta, positioning} % Librerie TikZ
\usepackage{pgfplots}         % Grafici avanzati
\pgfplotsset{compat=1.18}     % Impostazioni per pgfplots

% Pacchetti Matematici
\usepackage{amsmath, amssymb, amsthm} % Formule matematiche e teoremi

% Pacchetti per Tabelle e Unità di Misura
\usepackage{geometry}         % Margini della pagina
\usepackage{siunitx}          % Formattazione unità di misura
\usepackage{booktabs}         % Formattazione tabelle
\usepackage{tabularx}         % Tabelle avanzate
\usepackage{colortbl}         % Tabelle colorate

% Pacchetti per Layout e Formattazione
\usepackage{fancyhdr}         % Intestazione dei capitoli
\usepackage{float}            % Posizionamento di figure e tabelle
\usepackage{microtype}        % Migliora riempimento del testo

% Altri Pacchetti Utili
\usepackage{mdframed}         % Cornici
\usepackage{xcolor}           % Colore
\usepackage{animate}          % Animazioni

% Hyperref (da includere per ultimo)
\usepackage{hyperref}         % Link cliccabili


\definecolor{lightblue}{HTML}{9cd8fc} % Esempio di blu chiaro
% Definizione dei colori
\definecolor{myblue}{HTML}{139FE9}

\newmdenv[
    backgroundcolor=orange!10, % Scegli il colore e l'opacità dello sfondo
    linecolor=myblue,
    linewidth=0pt,
    leftmargin=0pt,
    rightmargin=0pt,
    innerleftmargin=10pt,
    innerrightmargin=10pt,
    innertopmargin=10pt,
    innerbottommargin=10pt
]{curiosita}

% Definizione dell'ambiente per i teoremi
\newmdtheoremenv[
  backgroundcolor=lightblue,
  linecolor=myblue,
  linewidth=1.5pt,
  innertopmargin=10pt,
  innerbottommargin=10pt,
  innerrightmargin=10pt,
  innerleftmargin=10pt
]{theorem}{Teorema}

% Definizione dell'ambiente per le definizioni
\newmdtheoremenv[
  backgroundcolor=gray!10,
  linewidth=0pt,
  innertopmargin=10pt,
  innerbottommargin=10pt,
  innerrightmargin=10pt,
  innerleftmargin=10pt
]{definition}{Definizione}

% Definizione dell'ambiente per le proposizioni
\newmdtheoremenv[
  backgroundcolor=gray!10,
  linecolor=myblue,
  linewidth=1.5pt,
  innertopmargin=10pt,
  innerbottommargin=10pt,
  innerrightmargin=10pt,
  innerleftmargin=10pt
]{proposition}{Proposizione}

% Configurazioni di fancyhdr
\pagestyle{fancy}
\fancyhf{} % Pulisce tutti i campi
\fancyhead[LE,RO]{\thepage} % Numero di pagina
\fancyhead[RE]{\nouppercase{\leftmark}} % Titolo del capitolo
\fancyhead[LO]{\nouppercase{\rightmark}} % Titolo della sezione
\renewcommand{\headrulewidth}{0.5pt}
\renewcommand{\footrulewidth}{0pt}

\geometry{a4paper, left=20mm, right=20mm, top=20mm, bottom=20mm}

\title{Esercizi Analisi Matematica T-A}
\author{Michele Mazza}
\date{November 2023}

\begin{document}

\frontmatter
\maketitle
\tableofcontents

\chapter*{Prefazione}
Il presente libro rappresenta una raccolta degli esercizi proposti nel corso di Analisi Matematica T-A tenuto dal Prof.re Cataldo Grammatico presso l'Università di Bologna. Rivolto agli studenti del corso di laurea in Ingegneria Gestionale (codice laurea 0925, Tipo di laurea L-9, Codice corso 28622).

\mainmatter
\chapter{Limiti}
\section{Limiti Base}

\subsection{Esercizio 01}

Dato il limite:

\[
    \lim_{{x \to -2}} \frac{{x^2 + x - 2}}{{x^2 + 5x + 6}}
\]

e l'equazione:

\[
    \lim_{{x \to -2}} \frac{{x^2 + x - 2}}{{x^2 + 5x + 6}} = -3
\]

Trova il seguente limite:

\[
    \lim_{{x \to -2}} \frac{{x^2 + x - 2}}{{x^2 + 5x + 6}}
\]

L'equazione si semplifica come segue:

\[
    \frac{{x^2 + x - 2}}{{x^2 + 5x + 6}} = \frac{{(x - 1)(x + 2)}}{{(x + 2)(x + 3)}}
\]

Quindi, il limite diventa:

\[
    \lim_{{x \to -2}} \frac{{(x - 1)(x + 2)}}{{(x + 2)(x + 3)}}
\]

Semplificando ulteriormente:

\[
    \frac{{(x - 1)(x + 2)}}{{(x + 2)(x + 3)}} = \frac{{x - 1}}{{x + 3}}
\]

Quindi, il limite diventa:

\[
    \lim_{{x \to -2}} \frac{{x - 1}}{{x + 3}}
\]

Poiché \(x - 1/x + 3\) è continuo in \(x = -2\), il limite diventa il valore di \(x - 1/x + 3\) in \(x = -2\).

Risposta:

\[
    \begin{aligned}
        \lim_{{x \to -2}} \frac{{x - 1}}{{x + 3}} & = -3
    \end{aligned}
\]
\subsection{Esercizio 02}

Dato il limite:

\[
    \lim_{{x \to a}} \frac{{\frac{1}{x} - \frac{1}{a}}}{{x - a}}
\]

e l'equazione:

\[
    \lim_{{x \to a}} \frac{{\frac{1}{x} - \frac{1}{a}}}{{x - a}} = -\frac{1}{a^2}
\]

Possibili passaggi intermedi:

Trova il seguente limite:

\[
    \lim_{{x \to a}} \frac{\frac{1}{x} - \frac{1}{a}}{x - a}
\]

L'equazione si semplifica come segue:

\[
    \frac{\frac{1}{x} - \frac{1}{a}}{x - a} = \frac{1}{x} - \frac{1}{a} \cdot \frac{1}{x - a}
\]

Quindi, il limite diventa:

\[
    \lim_{{x \to a}} \frac{1}{x} - \frac{1}{a} \cdot \frac{1}{x - a}
\]

Semplificando ulteriormente:

\[
    \frac{1}{x} - \frac{1}{a} \cdot \frac{1}{x - a} = \frac{1}{x} - \frac{1}{a} \cdot \frac{-a x}{1 - \frac{a}{x}}
\]

Quindi, il limite diventa:

\[
    \lim_{{x \to a}} \frac{1}{x} - \frac{1}{a} \cdot \frac{-a x}{1 - \frac{a}{x}}
\]

Semplificando ulteriormente:

\[
    \frac{1}{x} - \frac{1}{a} \cdot \frac{-a x}{1 - \frac{a}{x}} = -\frac{1}{a x}
\]

Quindi, il limite diventa:

\[
    \lim_{{x \to a}} -\frac{1}{a x}
\]

Poiché \(-\frac{1}{a x}\) è continuo in \(x = a\), il limite diventa il valore di \(-\frac{1}{a x}\) in \(x = a\).

Risposta:

\[
    \begin{aligned}
        \lim_{{x \to a}} -\frac{1}{a x} & = -\frac{1}{a^2}
    \end{aligned}
\]

\subsection{Esercizio 03}

Dato il limite:

\[
    \lim_{{x \to 4}} \frac{{\sqrt{x} - 2}}{{x^2 - 16}}
\]

e l'equazione:

\[
    \lim_{{x \to 4}} \frac{{\sqrt{x} - 2}}{{x^2 - 16}} = \frac{1}{32}
\]

Possibili passaggi intermedi:

Trova il seguente limite:

\[
    \lim_{{x \to 4}} \frac{{\sqrt{x} - 2}}{{x^2 - 16}}
\]

L'equazione si semplifica come segue:

\[
    \frac{{\sqrt{x} - 2}}{{x^2 - 16}} = \frac{{(\sqrt{x} - 2)(\sqrt{x} + 2)}}{{(x^2 - 16)(\sqrt{x} + 2)}} = \frac{1}{{(\sqrt{x} + 4)(\sqrt{x} + 2)}}
\]

Quindi, il limite diventa:

\[
    \lim_{{x \to 4}} \frac{1}{{(\sqrt{x} + 4)(\sqrt{x} + 2)}}
\]
\begin{center}
    \begin{tikzpicture}
        \begin{axis}[
                xlabel={$x$},
                ylabel={$y$},
                title={Grafico di $\frac{\sqrt{x} - 2}{x^2 - 16}$},
                xmin=0, xmax=4,
                ymin=0, ymax=0.25,
                grid=both,
                axis lines=middle,
                axis line style=->,
                grid style={line width=.1pt, draw=gray!10},
                major grid style={line width=.2pt,draw=gray!50},
                minor tick num=5,
                enlargelimits={abs=0.5},
                ticklabel style={font=\small},
                xlabel style={at={(ticklabel* cs:1)}, anchor=north west},
                ylabel style={at={(ticklabel* cs:1)}, anchor=south west},
                domain=(0:6),
                unbounded coords=jump, % Gestisce i punti non definiti
                samples=100, % Numero di campioni per la curva
                legend pos=south east
            ]
            \addplot[myblue, thick, smooth, mark=none] {1/((sqrt(x) + 4)*(sqrt(x) + 2))};
            \addplot[red, mark=*] coordinates{(4, 1/32)}; % Punto f(4) in rosso
            \draw[dashed] (4, 0) -- (4, 0.03); % Linea verticale in x=4
            \legend{{$\frac{\sqrt{x} - 2}{x^2 - 16}$}}
        \end{axis}
    \end{tikzpicture}
\end{center}

Poiché $\frac{1}{{(\sqrt{x} + 4)(\sqrt{x} + 2)}}$ è continua in $x = 4$, il limite diventa il valore di $\frac{1}{{(\sqrt{x} + 4)(\sqrt{x} + 2)}}$ in $x = 4$.

Risposta:

\[
    \begin{aligned}
        \lim_{{x \to 4}} \frac{1}{{(\sqrt{x} + 4)(\sqrt{x} + 2)}} & = \frac{1}{32}
    \end{aligned}
\]

\end{document}
